\vspace{1cm}
در این بخش، ابتدا به بررسی و تحلیل انواع ساختارهای طراحی گیت‌های جمع‌کننده‌ی کامل
\efnote{Full Adder}
پرداخته و عملکرد آن‌ها را از نظر تأخیر زمانی مورد مقایسه قرار می‌دهیم.
با توجه به حجم بالای شبیه‌سازی‌ها، ارزیابی مدارها در این بخش صرفاً در گوشه‌ی فرآیندی
TT
انجام شده است. همچنین دمای کاری برابر با
$25^\circ\text{C}$
و نسبت طراحی مناسب برای ترانزیستورها
($\beta$)
برابر با
$2$
فرض شده است.
\\
جهت ارزیابی دقیق، عملکرد منطقی و رفتار زمانی مدارها با استفاده از دو روش شبیه‌سازی مجزا مورد بررسی قرار می‌گیرد:
\begin{itemize}
    \item \textbf{شبیه‌سازی عملکردی:} برای تأیید صحت عملکرد منطقی مدار، از گام‌های زمانی (Step) بزرگ و فایل‌های (\texttt{.vec}) استفاده می‌شود. خروجی این شبیه‌سازی یک فایل با پسوند \texttt{.errX} خواهد بود که در صورت بروز هرگونه مغایرت میان خروجی شبیه‌سازی‌شده و رفتار مورد انتظار، خطاها را ثبت (Log) می‌کند. این شبیه‌سازی‌ها در پوشه‌های به نام \texttt{Func} قرار دارند.
    \item \textbf{شبیه‌سازی زمانی:} برای استخراج دقیق زمان تأخیر انتشار و تحلیل رفتار گذرا استفاده می‌شود. این شبیه‌سازی ها در پوشه‌های \texttt{Delay} قرار دارند.
\end{itemize}
خروجی رقم نقلی خروجی
($C_{out}$)
زمانی خروجی منطقی یک می‌دهد که حداقل دو ورودی از میان سه ورودی
$A$،
$B$
و
$C_{in}$
یک باشند. این رفتار منطقی بیانگر «تابع اکثریت»
\efnote{Majority}
است که رابطه‌ی ریاضی آن به صورت زیر تعریف می‌شود:
$$C_{out} = MAJ(A, B, C_{in}) = A B + B C_{in} + A C_{in}$$
همچنین خروجی حاصل‌جمع
($S$)
از تابع
XOR
روی سه ورودی تبعیت می‌کند که رابطه‌ی آن به صورت زیر است:
$$S = A \oplus B \oplus C_{in}$$
\vspace{0.65cm}
\newimage[0.5\textwidth]{images/FA-block.jpg}{
    شماتیک بلوک جمع‌کننده کامل تک‌بیتی
}{fig:fa-block}
\vspace{1cm}
طبق خواسته پروژه می‌بایست با $L=90nm$ و $W_{min}=90nm$ کار می‌کردیم ولی به دلیل محدودیت‌های کتابخانه از ابعاد پروژه اول (یعنی بررسی و طراحی گیت‌های پایه)استفاده کرده‌ایم: $L=180nm$ و $W_{min}=1\mu m$

\newcode[1]{spice}{codes/param.sp}{پارامتر‌های استفاده شده در شبیه‌سازی}{code:sim-param}

\newpage

\newcode[1]{spice}{codes/fa-sim.sp}{شبیه‌سازی HSPICE مربوط به عملکرد گیت‌های FA}{code:fa-sim}

\newcode[1]{spice}{codes/fa-func.sp}{ورودی شبیه‌سازی عملکرد مدارات جمع‌کننده کامل (تمام این شبیه‌سازی‌ها در دمای 25 درجه و در گوشه TT طراحی هستند)}{code:fa-func}

%%%%%%%%%%%%%%%%%%%%%%%%%%%%%%%%%%%%%%%%%%%%%%%%%%%%%%%%%%%%%%%%%%%%%%%%%%%%%%%%%%%%%%%%%%%%%%%%%%%%%%%%%%%%%%
%%%%%%%%%%%%%%%%%%%%%%%%%%%%%%%%%%%%%%%%%%%%%%%%%%%%%%%%%%%%%%%%%%%%%%%%%%%%%%%%%%%%%%%%%%%%%%%%%%%%%%%%%%%%%%
%%%%%%%%%%%%%%%%%%%%%%%%%%%%%%%%%%%%%%%%%%%%%%%%%%%%%%%%%%%%%%%%%%%%%%%%%%%%%%%%%%%%%%%%%%%%%%%%%%%%%%%%%%%%%%

\newpage

\subsection{
    طراحی
    \lr{Gate Level}
}
در این طراحی گیت جمع‌کننده خود را با استفاده از گیت‌های پایه‌ی دیگر همچون \lr{XOR} یا \lr{NAND} طراحی می‌کنیم.
این طراحی براساس تابع منطقی خروجی‌ها است که در سطح گیت‌ها مدل شده‌اند.

\newimage[0.5\textwidth]{images/FA.jpeg}{
    شماتیک بلوک جمع‌کننده کامل تک‌بیتی با ساختار سطح گیت
}{fig:fa-gate}

\newcode[1]{spice}{codes/fa.sp}{مدل SPICE گیت جمع‌کننده کامل}{code:fa}

\newimage[\textwidth]{images/sim/FA-func.png}{
    نتیجه شبیه‌سازی جمع‌کننده کامل
}{fig:fa-gate-func}

%%%%%%%%%%%%%%%%%%%%%%%%%%%%%%%%%%%%%%%%%%%%%%%%%%%%%%%%%%%%%%%%%%%%%%%%%%%%%%%%%%%%%%%%%%%%%%%%%%%%%%%%%%%%%%
%%%%%%%%%%%%%%%%%%%%%%%%%%%%%%%%%%%%%%%%%%%%%%%%%%%%%%%%%%%%%%%%%%%%%%%%%%%%%%%%%%%%%%%%%%%%%%%%%%%%%%%%%%%%%%
%%%%%%%%%%%%%%%%%%%%%%%%%%%%%%%%%%%%%%%%%%%%%%%%%%%%%%%%%%%%%%%%%%%%%%%%%%%%%%%%%%%%%%%%%%%%%%%%%%%%%%%%%%%%%%
\newpage

\subsection{
    طراحی
    \lr{NAND Based}
}
در این طراحی از سری کردن دو گیت \lr{XNAND} استفاده‌شده است.
این گیت یک ساختار دو ورودی و دو خروجی است که همزمان حاصل \lr{NAND} و  \lr{XOR} دو ورودی‌اش را خروجی می‌دهد و از 4 گیت \lr{NAND} تشکیل شده‌است.
ابتدا به بررسی این گیت و سپس به جمع‌کننده کامل می‌پردازیم.

\subsubsection*{
    گیت
    \lr{XNAND}
}

\newimage[0.5\textwidth]{images/XNAND.jpg}{
    شماتیک گیت XNAND با ساختار سطح گیت
}{fig:fa-gate}

\newcode[1]{spice}{codes/xnand.sp}{مدل SPICE گیت XNAND}{code:xnand}

با توجه با ساده‌سازی جبر بولی زیر می‌توان صحت مدار XNAND را بررسی کرد
$$(\; \overline{AB}\;)A + (\;\overline{AB}\;)B = (\bar{A}+\bar{B})A + (\bar{A}+\bar{B})B = \bar{A}A + \bar{B}A + \bar{A}B + \bar{B}B = \bar{B}A + \bar{A}B = A\oplus B$$
این طراحی در مواردی که تکنولوژی ما تنها شامل گیت NAND هست قابل پیاده‌سازی است و از جهت یکپارچه‌سازی مدار مورد استفاده قرار بگیرد.
در صفحه بعد شماتیک مدار را مشاهده می‌کنید. خروجی S از XOR شدن بیت نقلی ورودی و حاصل XOR دو بیت ورودی نتیجه می‌شود. همچنین محاسبۀ رقم نقلی خروجی که با تابع MAJ بدست می‌امد از رابطه بولی زیر محاسبه شده‌است(می‌دانیم که دو لایه NAND همان یک لایه \lr{AND-OR} هست):

$$
    \begin{aligned}
        AB + \text{Cin}(A\oplus B) & = AB + A\bar{B}\text{Cin} + B\bar{A}\text{Cin}                               \\
                                   & = AB + AB\text{Cin} + A\bar{B}\text{Cin} + B\bar{A}\text{Cin} + AB\text{Cin} \\
                                   & = AB + A\text{Cin}( B + \bar{B}) + B\text{Cin}( A + \bar{A})                 \\
                                   & = AB + A\text{Cin} + B\text{Cin} = \text{MAJ}(A,B,\text{Cin})
    \end{aligned}
$$


\newimage[0.6\textwidth]{images/FA-NAND.jpeg}{
    شماتیک مدار جمع‌کننده کامل تک‌بیتی \lr{NAND Based}
}{fig:fa-nand}

\newcode[1]{spice}{codes/fa-nand.sp}{مدل SPICE مدار جمع‌کننده کامل مبتنی بر گیت \lr{NAND2}}{code:fa-nand}

\newimage[\textwidth]{images/sim/FA-func-nand.png}{
    نتیجه شبیه‌سازی جمع‌کننده کامل
}{fig:fa-nand-func}


%%%%%%%%%%%%%%%%%%%%%%%%%%%%%%%%%%%%%%%%%%%%%%%%%%%%%%%%%%%%%%%%%%%%%%%%%%%%%%%%%%%%%%%%%%%%%%%%%%%%%%%%%%%%%%
%%%%%%%%%%%%%%%%%%%%%%%%%%%%%%%%%%%%%%%%%%%%%%%%%%%%%%%%%%%%%%%%%%%%%%%%%%%%%%%%%%%%%%%%%%%%%%%%%%%%%%%%%%%%%%
%%%%%%%%%%%%%%%%%%%%%%%%%%%%%%%%%%%%%%%%%%%%%%%%%%%%%%%%%%%%%%%%%%%%%%%%%%%%%%%%%%%%%%%%%%%%%%%%%%%%%%%%%%%%%%
\newpage

\subsection{
    طراحی آینه‌ای یا
    \lr{Compact Design}
}

قبل از پرداختن به این معماری، به توابع \lr{Self-dual} و همچنین سیگنال‌های معروفی که در ادامه در معماری‌های دیگر نیز باآن‌ها روبه‌رو می‌شویم می‌پردازیم.
\newimage[\textwidth]{images/FA-mirror.png}{
    شماتیک ترانزیستوری مدار جمع‌کننده تک بیت آینه‌ای
}{fig:fa-mirror}

رفتار مدار جمع‌کننده کامل (\lr{Full Adder}) نسبت به بیت‌های ورودی $A$ و $B$ را می‌توان به سه بخش عملکردی مجزا تقسیم‌بندی کرد:
\begin{itemize}
    \item \textbf{تولید (\lr{Generate}):}
          زمانی که هر دو ورودی $A$ و $B$ برابر با $1$ منطقی باشند، بدون توجه به مقدار بیت نقلی ورودی ($C_{in}$)، یک رقم نقلی خروجی تولید می‌شود ($C_{out} = 1$). رابطه منطقی آن به صورت زیر تعریف می‌شود:
          $$
              G = A B
          $$
    \item \textbf{انتشار (\lr{Propagate}):}
          زمانی که تنها یکی از دو ورودی $A$ یا $B$ برابر با $1$ منطقی باشد، رقم نقلی خروجی وابسته به بیت نقلی ورودی خواهد بود؛ به این معنی که بیت نقلی ورودی عیناً به خروجی منتقل یا منتشر می‌شود ($C_{out} = C$). این تابع منطقی را می‌توان به فرم زیر نشان داد:
          $$
              P = A \oplus B
          $$
          که در بهینه‌سازی های سطح ترانزیستور می‌توان آن را به صورت $A + B$ نیز مدل کرد.
    \item \textbf{کشتن (\lr{Kill}):}
          زمانی که هر دو ورودی $A$ و $B$ برابر با صفر منطقی باشند، رقم نقلی خروجی تحت هر شرایطی صفر خواهد شد ($C_{out} = 0$). به عبارت دیگر، رقم نقلی ورودی در این مرحله از بین می‌رود یا «کشته» می‌شود:
          $$
              K = \overline{A} \cdot \overline{B}
          $$
\end{itemize}

جدول درستی ارائه شده در شکل زیر، نشان می‌دهد که در هر سطر از جدول، تنها یکی از حالات بالا برقرار است (\lr{$G+P+K=1$}).
\newimage[0.6\textwidth]{images/FA-mirror-logic0.png}{
    جدول درستی مدار جمع‌کننده‌  کامل
}{fig:fa-truth-table}


یکی از ویژگی‌های ریاضی بسیار مهم جمع‌کننده کامل که پیاده‌سازی سخت‌افزاری آن را ساده‌تر می‌کند، خاصیت د‌وگان‌بودن
(\lr{Self-Dual}) توابع $S$ و $C_{out}$ است.
یک تابع منطقی $f(x_1, x_2, \dots, x_n)$ زمانی دوگان (\lr{Self-Dual}) نامیده می‌شود که متمم تابع با متمم کردن تک‌تک ورودی‌های آن برابر باشد:
$$
    \overline{f(A, B, C)} = f(\overline{A}, \overline{B}, \overline{C})
$$
توابع خروجی مجموع ($S$) و رقم نقلی خروجی ($C_{out}$) در جمع‌کننده کامل هر دو دارای این خاصیت فیزیکی-منطقی هستند. این ویژگی به اجازه می‌دهد ساختار ترانزیستوری نیمه بالایی (\lr{Pull-up}) و نیمه پایینی (\lr{Pull-down}) مدار را به صورت کاملاً متقارن نسبت به یکدیگر طراحی کنند.
با بهره‌گیری از خاصیت \lr{Self-Dual} و مفاهیم تولید، انتشار و کشتن کری ($G, P, K$)، می‌توان مدار جمع‌کننده کامل را به روش آینه‌ای (\lr{Mirror Design}) پیاده‌سازی کرد.

بر اساس جدول درستی، رابطه خروجی متمم بیت نقلی ($\overline{C_{out}}$) که یک تابع اقلیت (\lr{Minority}) است، به صورت زیر بیان می‌شود:
$$
    \overline{C_{out}} = \overline{AB + C(A+B)}
$$
پس از آن می‌توان از این سیگنال برای محاسبه حاصل بیت جمع استفاده کرد به صورتی که در شکل زیر قابل مشاهده است:
$$
    S = ABC + (A + B + C) \cdot \overline{C_{out}}
$$
به این ترتیب تعداد ترانزیستور کمتری مصرف می‌شود و انتظار داریم که تاخیر مداری بهبود پیدا کند.

\newimage[\textwidth]{images/FA-mirror-logic.png}{
    جدول درستی مدار جمع‌کننده‌ کامل
}{fig:fa-truth-table-2}

\newpage

\newcode[0.65]{spice}{codes/fa-mirror.sp}{مدل SPICE مدار جمع‌کننده کامل compact شده}{code:fa-mirror}

\newimage[0.7\textwidth]{images/sim/FA-func-mirror.png}{
    نتیجه شبیه‌سازی جمع‌کننده کامل
}{fig:fa-mirror-func}


%%%%%%%%%%%%%%%%%%%%%%%%%%%%%%%%%%%%%%%%%%%%%%%%%%%%%%%%%%%%%%%%%%%%%%%%%%%%%%%%%%%%%%%%%%%%%%%%%%%%%%%%%%%%%%
%%%%%%%%%%%%%%%%%%%%%%%%%%%%%%%%%%%%%%%%%%%%%%%%%%%%%%%%%%%%%%%%%%%%%%%%%%%%%%%%%%%%%%%%%%%%%%%%%%%%%%%%%%%%%%
%%%%%%%%%%%%%%%%%%%%%%%%%%%%%%%%%%%%%%%%%%%%%%%%%%%%%%%%%%%%%%%%%%%%%%%%%%%%%%%%%%%%%%%%%%%%%%%%%%%%%%%%%%%%%%

\newpage

\subsection{
    طراحی
    \lr{TG Based}
}

در این معماری، به جای استفاده از گیت‌های منطقی از \lr{Transmission Gate} استفاده می‌شود که به عنوان مالتی‌پلکسر خواهند بود. این امر به کاهش چشمگیر تعداد ترانزیستورها، بهبود تاخیر مسیرهای بحرانی و دستیابی به نوسان ولتاژ کامل (\lr{Full Swing}) در خروجی کمک می‌کند.

\newimage[\textwidth]{images/FA-Transmission.png}{
    شماتیک ترانزیستوری و در سطح گیت مدار جمع‌کننده تک بیت مبتنی بر \lr{Transmission Gate}
}{fig:fa-tg}

ساختار منطقی این معماری بر پایه سیگنال انتشار (\lr{Propagate}) یعنی $P = A \oplus B$ بنا شده است. با توجه به شکل مدار، ابتدا سیگنال انتشار $P$ و متمم آن $\overline{P}$ تولید می‌شود. همان‌طور که در شماتیک بلوکی و ترانزیستوری مشخص است، این پیاده‌سازی را می‌توان به دو بخش اصلی تقسیم کرد:

پس از تولید سیگنال‌های کنترلی $P$ و $\overline{P}$، خروجی‌های اصلی با استفاده از ساختار مالتی‌پلکسری مبتنی بر کلیدهای انتقال پیاده‌سازی می‌شوند:
\begin{itemize}
    \item \textbf{محاسبه مجموع ($S$):}
          خروجی $S$ از رابطه منطقی $S = P \oplus C_{in}$ به دست می‌آید. برای پیاده‌سازی این رابطه، از دو کلید انتقال استفاده می‌کنیم که مثل یک کلید دوحالته رفتار می‌کنند؛ به این صورت که با توجه به وضعیت سیگنال‌های $P$ و $\overline{P}$، تصمیم می‌گیرند که کدام‌یک از سیگنال‌های $C_{in}$ یا متمم آن یعنی $\overline{C_{in}}$ به سمت خروجی هدایت شوند. در نهایت هم یک اینورتر در خروجی قرار داده‌ایم تا سیگنال خروجی $S$ تقویت شده و قدرت کافی برای تحریک مراحل بعدی را داشته باشد.
          $$
              \overline{S} = P \cdot C_{in} + \overline{P} \cdot \overline{C_{in}} \implies S = P \oplus C_{in}
          $$

    \item \textbf{محاسبه رقم نقلی خروجی ($C_{out}$):}
          تولید رقم نقلی مستقیماً به وضعیت سیگنال انتشار ($P$) بستگی دارد. اگر $P = 1$ باشد (یعنی فقط یکی از ورودی‌های $A$ یا $B$ یک باشند)، رقم نقلی خروجی دقیقاً همان بیت نقلی ورودی یعنی $C_{in}$ خواهد بود. اما اگر $P = 0$ باشد، مقدار خروجی دیگر به $C_{in}$ ربطی ندارد و مستقیماً از روی ورودی $A$ تعیین می‌شود. این مسیر هم در انتها به یک اینورتر ختم می‌شود تا سیگنال نهایی $C_{out}$ بدون افت ولتاژ و کاملاً تقویت‌شده به دست بیاید.
          $$
              \overline{C_{out}} = P \cdot \overline{C_{in}} + \overline{P} \cdot \overline{A} \implies C_{out} = P \cdot C_{in} + \overline{P} \cdot A
          $$
\end{itemize}

\newcode[0.78]{spice}{codes/fa-tg.sp}{مدل SPICE مدار جمع‌کننده کامل مبتنی بر TG}{code:fa-tg}

\newimage[0.8\textwidth]{images/sim/FA-func-tg.png}{
    نتیجه شبیه‌سازی جمع‌کننده کامل
}{fig:fa-tg-func}



% %%%%%%%%%%%%%%%%%%%%%%%%%%%%%%%%%%%%%%%%%%%%%%%%%%%%%%%%%%%%%%%%%%%%%%%%%%%

\newpage

\subsection{مقایسه انواع طراحی‌ها}

\indent
ازانجایی که مدارات ما در عملکرد، یکسان عمل می‌کنند و همگی درست هستند برای مقایسه آن‌ها باید از پارامتر تاخیر و توان مصرفی استفاده کرد.
در ادامه توان مصرفی مدار‌ها را، همچنین تاخیر انتشار بیت نقلی و بیت مجموع را با استفاده از شبیه‌سازی‌های جداگانه‌ای اندازه‌گیری کرده و خروجی‌ها را مقایسه می‌کنیم.
در انتها بهترین پیاده‌سازی را برای استفاده در جمع‌کننده‌های چند بیتی استفاده می‌کنیم.

ما از اسکریپت‌های پایتون \texttt{delay-sim.py} و \texttt{modify-mtx.py} برای استخراج دقیق تاخیرها و توان مصرفی از فایل‌های خروجی شبیه‌سازی استفاده کرده‌ایم. این اسکریپت‌ها به صورت خودکار داده‌های مربوط به تأخیر انتشار بیت نقلی ($C_{out}$) و بیت مجموع ($S$) را از فایل‌های \texttt{.mtx} استخراج کرده و در قالب فایل \texttt{.json} ذخیره می‌کنند. این فایل‌های \texttt{.json} سپس برای تجزیه و تحلیل بیشتر و مقایسه عملکرد طراحی‌های مختلف استفاده خواهند شد.
\pfnote{این اسکریپت‌ها توسط \lr{LLM-Tools} تولید شده اند ولی هزار بار چک شده اند :D}
با این روش در هر idex از شبیه‌سازی ما یکی از حالات transition را بررسی می‌کنیم و مقادیر تاخیر خود را اندازه‌گیری می‌کنیم.

\newimage[\textwidth]{images/FA-compare.png}{
    مقایسه انواع طراحی‌های مختلف جمع‌کننده کامل
}{fig:fa-compare}

\newimage[\textwidth]{images/FA-delay-compare.png}{
    مقایسه انواع تاخیر در طراحی‌های مختلف جمع‌کننده کامل
}{fig:fa-delay-compare}

همانطور که در نمودارهای بالا مشاهده می‌شود، طراحی مبتنی بر \lr{Transmission Gate} به دلیل استفاده بهینه از ترانزیستورها و ساختار مالتی‌پلکسری، عملکرد بهتری در مقایسه با سایر معماری‌ها ارائه می‌دهد. این طراحی نه تنها تاخیر کمتری دارد بلکه توان مصرفی آن نیز نسبت به سایر پیاده‌سازی‌ها کاهش یافته است. در نتیجه، این معماری برای استفاده در جمع‌کننده‌های چند بیتی بسیار مناسب‌تر است.
البته که همچنان هم در برخی موارد که تاخیر کمتری نیاز هست می‌توان از طراحی آینه‌ای استفاده کرد ولی در کل طراحی مبتنی بر \lr{Transmission Gate} به عنوان بهترین گزینه برای جمع‌کننده‌های چند بیتی انتخاب می‌شود.
همینطور از این نکته هم نمی‌شود غافل شد که طراحی \lr{NAND-based} به دلیل استفاده صرف از \lr{NAND2} در طراحی‌های اتوماتیک راحت‌تر مورد استفاده قرار می‌گیرد.
در ادامه تمام مدارات جمع‌کنند از طراحی مبتنی بر \lr{Transmission Gate} استفاده خواهند کرد.

\newimage[0.9\textwidth]{images/sim/FA-delay.png}{
    نتیجه شبیه‌سازی جمع‌کننده کامل در سطح گیت
}{fig:fa-delay}



\newcode[0.76]{spice}{codes/delay.sp}{نمونه کد HSPICE مورد استفاده برای اندازه‌گیری تاخیر}{code:delay}